
\subsubsection{Token-Probability Methods}

The simplest UQ approach for language models uses the model's predicted probability distribution directly. Negative log-probability and maximum softmax probability are standard baselines (Kadavath et al., 2022; Malinin and Gales, 2021). Despite their simplicity, these methods have demonstrated competitive performance across benchmarks. In the original SE paper, Farquhar et al. (2024) report token-probability AUROC of approximately 0.67 on TriviaQA for Llama-2-70B. The present study finds token-probability AUROC of 0.6835 on TriviaQA and 0.6551 on NaturalQuestions for Llama-3-8B-Base—values broadly consistent with these prior reports.

\subsubsection{Sampling-Based Semantic Methods}

Farquhar et al. (2024) introduced SE to address the sensitivity of logit-based estimates to paraphrase variation: two responses expressing the same answer with different wording should receive similar uncertainty scores, but token-probability can differ substantially. SE clusters N stochastic responses into semantic equivalence classes via bidirectional entailment checking (using DeBERTa-large-mnli) and computes entropy over the cluster probability mass. The method was evaluated on Llama-2 variants (multiple scales), GPT-3, and Falcon models. Reported AUROC values range from 0.72 to 0.79 on TriviaQA, NaturalQuestions, BioASQ, and SQuAD, consistently exceeding token-probability.

KLE (Nikitin et al., 2024) generalizes SE using von Neumann entropy over a semantic similarity kernel matrix, computed from pairwise NLI scores. It provides a theoretically unified framework but shares SE's dependence on semantic diversity across samples: when the similarity matrix has rank 1 (all samples identical), von Neumann entropy is near zero.

SEPs (Kossen et al., 2024) approximate SE from hidden-state representations using linear probes trained on a reference set, achieving 5–10× inference speedup relative to multi-sample SE. SEPs were evaluated on Llama-2 variants on TriviaQA and NaturalQuestions. The degenerate_fraction of the generation distribution is not reported in these evaluations.

SelfCheckGPT (Manakul et al., 2023) measures consistency across N samples using NLI, BERTScore, or n-gram overlap as consistency scoring functions. Evaluated primarily on GPT-3 outputs, SelfCheckGPT-NLI achieves strong performance on biographical fact-checking. Unlike SE and KLE, SelfCheckGPT-NLI does not require multiple distinct clusters; it measures pairwise entailment consistency, which can retain some discriminative signal even when most samples are near-identical, provided the minority of diverse samples carry entailment information.

\subsubsection{Output Diversity in Language Generation}

The relationship between sampling diversity and generation quality has been studied in other contexts. Diverse Beam Search (Vijayakumar et al., 2016) was proposed to address near-identical outputs under standard beam decoding. Self-BLEU (Zhu et al., 2018) and distinct-n (Li et al., 2016) measure corpus-level lexical diversity in generation. The \texttt{degenerate_fraction} metric introduced in this paper differs from these in focus: it measures the proportion of queries for which semantic clustering collapses to K=1, which is the condition under which SE and KLE become uninformative as UQ estimators. The application of diversity measurement as a pre-screening diagnostic for UQ validity is, to the authors' knowledge, not established in prior work.

\subsubsection{Representation-Based Methods}

Several methods bypass sampling by extracting uncertainty estimates from internal model representations. Layer-wise Semantic Dynamics (LSD; Mir, 2025) trains contrastive representations on hidden states, reporting AUROC of 0.96 on TruthfulQA. Effective Rank-based Uncertainty (Wang et al., 2025) uses spectral rank of hidden-state matrices. Conformal prediction methods (Cherian et al., 2024; Su et al., 2024) provide coverage guarantees. These approaches do not require sampling diversity and are thus not subject to the failure mode described here.

\subsubsection{Surveys}

Kang et al. (2025) survey over 100 UQ methods for LLMs and identify evaluation fragmentation—particularly the conflation of methods evaluated under different conditions—as a primary limitation of the field. The present finding that base and instruction-tuned models produce substantially different sampling distributions, and that this distinction is not consistently reported in UQ benchmarks, is one instance of this fragmentation.

